\chapter*{국문요약}
\addcontentsline{toc}{chapter}{국문요약}
\begin{center}
\Large\textbf{기하학적 경로계획, 강화학습, 힘 인지 모방학습 및 특징 마스킹 기반 하이브리드 로봇 표면 마감 프레임워크 개발}
\end{center}

본 학위논문 초안은 하나의 제품 전체를 자동으로 폴리싱하기 위한 하이브리드 로봇 표면 마감 프레임워크를 제안한다. 제안 프레임워크는 CAD 또는 재구성된 표면 전체를 커버하는 공구 폭 기반 전역 경로계획, 실제 접촉력에 따라 sliding velocity를 조절하는 강화학습, 에지 및 국소 난구간에서 위치·자세·목표 힘 궤적을 생성하는 힘 인지 모방학습, 그리고 반사광과 배경 정보에 대한 의존을 줄이는 특징 마스킹으로 구성된다.

전역 경로계획에서는 TCP 자체가 아닌 실제 polishing pad의 유효 접촉 폭과 overlap ratio를 이용해 경로 간격을 결정하고, 공구가 이동하며 형성하는 swept footprint의 합집합이 목표 표면을 포함하도록 full-coverage path를 생성한다. 강화학습 정책은 로봇 관절이나 6차원 자세를 직접 출력하지 않고, 고정된 기준 경로의 진행 속도 scale만 조절하여 Preston equation 기반 가공률 편차를 줄인다. 힘 인지 모방학습 정책은 RGB 영상, 로봇 상태 및 6축 힘·토크 이력을 입력으로 받아 위치 3축, 자세 3축, 목표 힘 3축으로 이루어진 9차원 action chunk를 생성하고, 저수준 admittance controller가 이를 추종한다.

특징 마스킹 장에서는 VIOLA의 object proposal 기반 object-centric representation과 Mask2Act의 future object-mask latent plan을 현재 IL framework에 적용하는 방법을 제안한다. 대상 surface, tool footprint, 보호해야 할 hardware 및 결함 영역을 mask로 표현하고, object-centric region token과 미래 mask sequence를 힘 정보 및 로봇 상태와 융합해 9차원 pose-force action을 생성한다. 이 부분은 아직 구현 전의 설계안으로서 정적 mask conditioning, VIOLA-inspired region token, Mask2Act-inspired future-mask plan 및 결합 모델을 비교하는 실험을 계획한다.

실증 대상물은 넓은 도장 곡면과 에지, 컷어웨이, 프렛, 픽업, 브릿지 등 다양한 국소 금속 구간이 하나의 제품에 공존하는 기타로 설정한다. 최종 연구는 전체 커버리지, 가공량 균일도, 접촉력 안정성, 안전성, 교시 및 프로그래밍 효율과 다양한 형상 및 사용 상태에 대한 일반화 성능을 종합적으로 평가한다.

\vspace{0.5cm}
\noindent\textbf{주요어:} 로봇 표면 마감, 전면 커버리지 경로계획, 공구 접촉 폭, 강화학습, 힘 인지 모방학습, 특징 마스킹, 어드미턴스 제어
